
\documentclass[12pt,a4paper]{article}

\title{Métodos Computacionais em Física de Altas Energias - Aula 1}
\author{Theo Basso Pilon (RA: 219520)}
\date{}

\begin{document}

\maketitle

\section{Introdução}

A aula desta semana apresentou ferramentas computacionais importantes para o desenvolvimento de software em Física de Altas Energias. Para a primeira atividade, vamos utilizar algumas ferramentas discutidas, sendo elas o \LaTeX, \texttt{make} e \texttt{sed}. Parte disso será feita através da criação e manutenção básica de um repositório no Github, introduzindo comandos básicos de git, bem como a escrita de um documento em \LaTeX, discutindo brevemente os tópicos abordados nessa primeira aula. O curso também irá abordar algumas ferramentas de análise de dados, como o ROOT \cite{Brun:1997pa}, por exemplo. 


\section{Conceitos vistos nos slides}

Dentre os conceitos vistos na aula, vamos definir três principais:

\begin{itemize}
    \item \textbf{Git:} Criação de um repositório para um projeto, que pode ser compartilhado e editado por várias pessoas. Permite acompanhar as alterações dos arquivos, criar commits , trabalhar com branches e sincronizar o projeto com repositórios remotos.
    \item \textbf{Make:} Determina quais partes de um programa precisam ser recompiladas e executa os comandos necessários. As regras de um projeto são organizadas em um arquivo \texttt{Makefile}, relacionando targets, dependências e comandos.
    \item \textbf{Sed:} editor de fluxo (stream editor) utilizado para realizar transformações básicas em texto. Um exemplo apresentado nos slides é a substituição de uma string por outra usando \texttt{sed 's/hello/world/' input.txt}.
\end{itemize}


\section{Conclusão}

Git, make e sed são ferramentas complementares e fundamentais em projetos computacionais. O Git auxilia no controle e na organização das versões do código, o make automatiza etapas de compilação e o sed permite realizar transformações de texto de maneira eficiente. Juntamente com \LaTeX{}, essas ferramentas serão bastante usadas durante o curso. Posteriormente, vamos abordar questões mais complexas envolvendo análise de dados usando o ROOT como ferramenta principal.

\bibliographystyle{unsrt}
\bibliography{referencias}
\end{document}

